% title:          Continuous function meeting a derivative
% area:           limits-continuity, derivatives
% methods:        darboux-theorem, intermediate-value-theorem
% difficulty:     easy
% difficulty-ai:  easy
% status:
% review:         human
% --- history: all optional, leave EMPTY if unknown ---
% origin:         imc
% origin-date:    2019-07-31
% origin-number:  Day 2, Problem 6
% also-in:
% references:     Earliest known appearance (to the best of our knowledge): IMC 2019, Day 2, Problem 6 - https://www.imc-math.org.uk/imc2019/imc2019-day2-questions.pdf
% history:        ai
\begin{problem}
    Let \( f, g \in C^0\left(\mathbb{R},\mathbb{R}\right) \) be such that \( g \) is differentiable
    and
    \[
        \left(f\left(0\right)-g'\left(0\right)\right)\left(g'\left(1\right)-f\left(1\right)\right) >
        0.
    \]
    Show that there exists \( c \in \left]0;1\right[ \) with \( f\left(c\right) =
    g'\left(c\right) \).
\end{problem}
\begin{solution}
    Let \( F := f-g' \colon \mathbb{R} \to \mathbb{R} \). Then the condition reads
    \[
        0 <
        \left(f\left(0\right)-g'\left(0\right)\right)\left(g'\left(1\right)-f\left(1\right)\right) =
        -F\left(0\right)F\left(1\right), \quad \text{i.e.} \quad F\left(0\right)F\left(1\right) < 0.
    \]
    In particular, \( F\left(0\right), F\left(1\right) \neq 0 \) and
    \( -\operatorname{sign}\left(F\left(0\right)\right) =
    \operatorname{sign}\left(F\left(1\right)\right) \).
    \\\\
    If \( g' \) were continuous, then \( F \) would be continuous, so that, by the intermediate
    value theorem, \( F\left[\left[0;1\right]\right] \) would be an interval containing
    \( F\left(0\right) \) and \( F\left(1\right) \), which have opposite signs. Hence, \( 0 \in
    F\left[\left[0;1\right]\right] \), which means that there would exist \( c \in
    \left[0;1\right] \) with \( F\left(c\right) = 0 \), i.e., \( f\left(c\right) =
    g'\left(c\right) \), and clearly \( c \in \left]0;1\right[ \) as \( F\left(0\right),
    F\left(1\right) \neq 0 \).
    \\\\
    Now, \( g' \) is not necessarily continuous, but we can bypass this. Define
    \[
        H \colon \mathbb{R} \to \mathbb{R}, \quad x \mapsto \int_0^x f\left(t\right) \,
        \mathrm{d}t-\left(g\left(x\right)-g\left(0\right)\right),
    \]
    which is well defined as \( f \) is continuous. By the fundamental theorem of calculus applied
    to \( f \) and since \( g \) is differentiable, \( H \) is differentiable with \( H' = f-g' =
    F \), so that \( H'\left(0\right)H'\left(1\right) < 0 \). Then, since \( H \) is differentiable
    on the interval \( \mathbb{R} \), \( H' \) has the Darboux property by Darboux's
    theorem [\appendixref{darboux-theorem}], so that \( H'\left[\left[0;1\right]\right] \) is an
    interval containing \( F\left(0\right) \) and \( F\left(1\right) \), which have opposite signs.
    Hence, \( 0 \in H'\left[\left[0;1\right]\right] \), which means that there exists \( c \in
    \left[0;1\right] \) with \( H'\left(c\right) = 0 \), i.e., \( f\left(c\right) =
    g'\left(c\right) \), and clearly \( c \in \left]0;1\right[ \) as \( F\left(0\right),
    F\left(1\right) \neq 0 \).
\end{solution}
